\documentclass[11pt]{article}
\usepackage{graphicx} % Required for inserting images
\setlength{\parindent}{0pt}
\usepackage{hyperref}
\usepackage{enumitem}
\usepackage[utf8]{inputenc} 
\usepackage[T1]{fontenc}
\usepackage[brazil]{babel}
\usepackage{lipsum}
\usepackage[left=1.06cm,top=1.7cm,right=1.06cm,bottom=0.49cm]{geometry}

%by: Aline R. Antunes

\begin{document}
	\begin{center}
		\textbf{Kaleb Irahola Azad}\\ 
		\hrulefill
	\end{center}
	
	\begin{center}
		Calle {Álvarez 1822} \textbullet\ Viña del mar, Chile \textbullet\ kaleb.irahola.a@mail.pucv.cl \textbullet\ +56 950532667
	\end{center}
	
	\vspace{0.5pt}
	
	Profesional en Ingeniería Mecatrónica enfocado en el diseño y gestión de proyectos tecnológicos que conectan la universidad con el sector productivo y el entorno local. Impulsa la robótica, la automatización y la I+D+i aplicada como herramientas para formar talento crítico y generar soluciones concretas a problemas reales.
	
	\vspace{0.5pt}
	
	\begin{center}
		\textbf{Experiencia}
	\end{center}
	\textbf{Pontificia Universidad Católica de Valparaíso} \hfill Valparaíso, Chile
	
	\textbf{Doctorante - Escuela de Ingeniería Industrial} \hfill Marzo de 2026 – Actualidad
	\begin{itemize}[noitemsep, topsep=0pt, partopsep=0pt, parsep=0pt]
		\item Doctorante e Investigador en la Escuela de Ingeniería Industrial
	\end{itemize}
	
	\textbf{Universidad Católica Boliviana "San Pablo"} \hfill Tarija, Bolivia
	
	\textbf{Docente Tiempo Completo - Carrera de Ingeniería Mecatrónica} \hfill Febrero de 2022 – Enero de 2026
	\begin{itemize}[noitemsep, topsep=0pt, partopsep=0pt, parsep=0pt]
		\item Impartir 34 asignaturas semestrales y 2 asignaturas modulares de pregrado en Ingeniería Mecatrónica, abarcando áreas como robótica, automatización industrial, instrumentación, electrónica, matemáticas y gestión de proyectos.
		\item Tutorizar, relatar y monitorear trabajos de grado de 33 estudiantes de la Carrera Ingeniería Mecatrónica.
	\end{itemize}
	
	\vspace{12pt}
	
	\begin{center}
		\textbf{Liderazgo y Actividades}
	\end{center}
	
	\textbf{Universidad Católica Boliviana "San Pablo"}	\hfill Tarija, Bolivia
	
	\textbf{Docente Tiempo Completo - Carrera de Ingeniería Mecatrónica} \hfill Febrero de 2022 – Enero de 2026
	\begin{itemize}[noitemsep, topsep=0pt, partopsep=0pt, parsep=0pt]
		\item Líder del Re-diseño curricular de la Carrera de Ingeniería Mecatrónica para elaborar el Plan de Estudios 2025.
		\item Líder del Comité de Autoevaluación con fines a la Acreditación de la Carrera de Ingeniería Mecatrónica ante el Comité Ejecutivo de Universidades Bolivianas (CEUB) y la ejecución del plan de mejoras.
		\item Encargado del Re-diseño y actualización de la infraestructura y guías de uso de los laboratorios de Automatización y Robótica, estandarizando prácticas, protocolos de seguridad y criterios de evaluación.
		\item Coordinador del proyecto: Traductor de lengua de señas boliviana a comunicación Oral.
		\item Colaborar en el proyecto: Ecotintes, para la revalorización de la cultura misio-guaraní en la comunidad de Chimeo en la aplicación de textiles.
		\item Impulsar la transferencia de proyectos de grado al sector productivo y a semilleros municipales, gestionando la obtención de capital semilla y la articulación con actores locales.
	\end{itemize}
	
	\vspace{12pt}
	
	\textbf{Comité académico nacional de la OCEPB}	\hfill Tarija, Bolivia
	
	\textbf{Encargado Departamental – Área de Robótica} \hfill 2022 – 2025
	\begin{itemize}[noitemsep, topsep=0pt, partopsep=0pt, parsep=0pt]
		\item Organizador de la etapa departamental del área de robótica de la Olimpiada Científica Estudiante Plurinacional Boliviana, coordinando la capacitación de docentes y estudiantes, y liderando la evaluación técnica de los equipos.
	\end{itemize}
	
	\vspace{12pt}
	
	\textbf{Startup Weekend Bolivia}	\hfill El Alto, Bolivia
	
	\textbf{Ganador del Startup Weekend Bolivia "Tech"} \hfill 2016
	\begin{itemize}[noitemsep, topsep=0pt, partopsep=0pt, parsep=0pt]
		\item Ganador del Startup Weekend Bolivia y concursante en el Global StartUp Battle con el proyecto AlertaUV.
	\end{itemize}
	
	\vspace{12pt}
	
	\textbf{8th International Olympiad on Astronomy and Astrophysics}	\hfill Suceava, Rumania
	
	\textbf{Representante de Bolivia} \hfill 2014
	\begin{itemize}[noitemsep, topsep=0pt, partopsep=0pt, parsep=0pt]
		\item Único representante de Bolivia en las 8vas Olimpiadas Internacionales de Astronomía y Astrofísica.
	\end{itemize}
	
	\newpage
	
	
	\begin{center}
		\textbf{Formación académica}
	\end{center}
	\textbf{Universidad Internacional de La Rioja} \hfill CDMX, México
	
	Maestro en Diseño y Gestión de Proyectos Tecnológicos. GPA: 8.78/10.0 \hfill 31 de julio de 2024
	
	\vspace{12pt}
	
	\textbf{Universidad Católica Boliviana "San Pablo"} \hfill Tarija, Bolivia
	
	Licenciado en Ingeniería Mecatrónica. Graduado por Excelencia, GPA: 92.94/100.0 \hfill 25 de septiembre de 2020
	
	\begin{center}
		\textbf{Formación complementaria}
	\end{center}
	
	\textbf{Universidad Católica Boliviana "San Pablo"} \hfill (En línea) Tarija, Bolivia
	
	Diplomado en Análisis de Datos con Inteligencia Artificial. \hfill 2026
	
	\vspace{12pt}
	
	\textbf{Universidad Amazónica de Pando} \hfill (En línea) Cobija, Bolivia
	
	Diplomado en Educación Superior. \hfill Abril de 2024
	
	\vspace{12pt}
	
	\textbf{Universidad Internacional de La Rioja} \hfill (En línea) CDMX, México
	
	Diplomado en Auditoría, Evaluación y Legislación de Proyectos Tecnológicos. \hfill Noviembre de 2023
	
	\vspace{12pt}
	
	\textbf{Universidad Internacional de La Rioja} \hfill (En línea) CDMX, México
	
	Diplomado en Dirección de Proyectos. \hfill Mayo de 2023
	
	\vspace{12pt}
	
	\textbf{Universidad Internacional de La Rioja} \hfill (En línea) CDMX, México
	
	Diplomado en Innovación Tecnológica. \hfill Noviembre de 2022
	
	\vspace{12pt}
	
	\textbf{Universidad Católica Boliviana "San Pablo"} \hfill (En línea) Tarija, Bolivia
	
	Curso de Investigación Científica con Inteligencia Artificial. \hfill Julio de 2025
	
%	\vspace{12pt}
%	
%	\textbf{Universidad Mayor de San Andrés} \hfill (En línea) La Paz, Bolivia
%	
%	XXII Curso Boliviano de Sistemas Complejos. \hfill Noviembre de 2024
	
	\vspace{12pt}
	
	\textbf{Universidad Católica Boliviana "San Pablo"} \hfill (Híbrido) La Paz, Bolivia
	
	Curso de Evaluación y Acreditación de la Calidad de Programas Académicos. \hfill Julio de 2023
	
	\vspace{12pt}
	
	\textbf{Universidad Católica Boliviana "San Pablo"} \hfill (En línea) Tarija, Bolivia
	
	Curso de Machine Learning con Python. \hfill Junio de 2023
	
	\vspace{12pt}
	
	\textbf{Ideas Automatitation} \hfill (En línea) Cochabamba, Bolivia
	
	Curso de Especialidad en Automatización y Control con Schneider Electric. \hfill Abril de 2020
		
	\begin{center}
		\textbf{Habilidades e Intereses}
	\end{center}
	
	\textbf{Técnicas:} Python (entorno de análisis), ROS~2, Linux/Ubuntu, LaTeX, Git y GitHub; Métodos cuantitativos; Diseño y simulación CAD, prototipado rápido; Automatización e instrumentación industrial, PLCs, sistemas HMI, redes industriales básicas; Electrónica analógica y digital; Sistemas embebidos y microcontroladores; Microsfot 365.
	
	\vspace{4pt}
	
	\textbf{Idiomas:} Español (nativo); Inglés (básico); Portugués (básico).
	
	\vspace{4pt}
	
	\textbf{Intereses:} Proyectos tecnológicos para el desarrollo regional; Sistemas Complejos; Modelos Matemáticos; Análisis y optimización de sistemas; Educación en ingeniería y aprendizaje activo; actividades de divulgación científica.
	
	
\end{document}
